\section{Introdução}

O primeiro modelo atômico foi idealizado por J. Thomson, por volta de 1910. 
De acordo com ele, os elétrons eram cargas negativas puntiformes imersas em uma distribuição contínua 
e esféricamente simétrica de carga positiva. 
Os elétrons em um átomo excitado vibrariam em torno de suas posições de equilíbrio, explicando 
a emissão de radiação. No entanto, os espectros observados experimentalmente iam de encontro ao modelo de Thomson.

Foi um ex-aluno de Thomson, E. Rutherford, que refutou de maneira definitiva esse modelo. 
Realizando experimentos sobre o espalhamento de partículas $\alpha$ (positivamente carregadas) por átomos de ouro, 
Rutherford concluiu que a toda a carga positiva do átomo estava concentrada em uma pequena região em seu centro, que era ordens de grandeza menor que o valor do raio atômico. 
Essa região recebeu o nome de núcleo atômico, e possui praticamente toda a massa de um átomo. 

O modelo de Rutherford, no entanto, nada falava sobre a distribuição dos elétrons dentro do átomo. 
Por volta de 1913, N. Bohr, que trabalhava no laboratório de Rutherford na época dos experimentos de Geiger e Mardsen, formulou os seguintes postulados sobre a estrutura atômica:

\begin{itemize}
\item[1.] Um elétron em um átomo se move em uma órbita circular em torno do núcleo sob influência 
da atração coulombiana entre o elétron e o núcleo.

\item[2.] Em vez da infinidade de órbitas que seriam possíveis segundo a mecânica clássica, um elétron
só pode se mover em uma órbita na qual seu momento angular orbital L é um múltiplo
inteiro de $\hbar $ (a constante de Planck dividida por $2\pi$).

\item[3.] Apesar de estar constantemente acelerado, um elétron que se move em uma dessas órbitas
possíveis não emite radiação eletromagnética. Portanto sua energia total E permanece constante.

\item[4.] É emitida radiação eletromagńetica se um elétron, que se move inicialmenle sobre uma órbita 
de energia total $E_i$, muda seu movimenlo descontinuamente de forma a se mover em uma
órbita de energía total $E_f$. A freqüência da radiação emitida $\nu$ é igual à quantidade $(E_i - E_f)$
dividida pela constante de Planck $h$.
\end{itemize}

Embora modelos mais precisos do que o de Bohr tenham sido desenvolvidos posteriormente, ele serve bem para
realizar previsões teóricas que concordam com os dados experimentais. 
Em especial, os postulados 2, 3, e 4 dizem que a energia é absorvida e emitida de forma quantizada pelo átomo,
fato este que foi experimentalmente comprovado em 1914 por J. Franck e G. Hertz em um simples experimento, que será o 
tópico deste trabalho. \cite{eisberg} \cite{tipler}